%==============================================================================
%  Starting a PhD in ML / NLP -- a practical starter guide
%  Konstantinos Skianis, University of Ioannina
%  Compile:  latexmk -cd -pdf phd-ml-nlp-starter-guide.tex
%==============================================================================
\documentclass[11pt,a4paper]{article}
\usepackage{guide-style}
\guidemeta{Starting a PhD in ML / NLP}{Starting a PhD in Machine Learning or NLP: A Practical Starter Guide}{PhD onboarding, ML/NLP research practice}

\begin{document}

%==============================================================================
\begin{titlepage}
\centering
\vspace*{2.2cm}
{\color{brandlight}\rule{\textwidth}{2pt}}\\[0.8cm]
{\Huge\bfseries\color{brand} Starting a PhD in\\[0.25cm] Machine Learning or NLP\par}
\vspace{0.7cm}
{\LARGE\color{ink} A Practical Starter Guide\par}
\vspace{0.5cm}
{\large\color{mutegray} Venues, deadlines, tools, and the habits that compound\par}
\vspace{0.6cm}
{\color{brandlight}\rule{\textwidth}{2pt}}\\[1.1cm]

{\large\bfseries\color{ink} Konstantinos Skianis\par}
\vspace{0.18cm}
{\normalsize\color{mutegray} Assistant Professor of Natural Language Processing\\[2pt]
Department of Computer Science and Engineering\\[2pt]
University of Ioannina\par}
\vspace{0.25cm}
{\footnotesize\color{brandlight}\url{https://y3nk0.github.io}\par}
\vspace{1.1cm}

\begin{minipage}{0.85\textwidth}
\small\color{ink}
Everything here is operational: what to set up in your first week, which venues
your work will target, how to find and judge related work, how to keep a
library that never loses a reference, how to use an AI assistant without
damaging your scholarship, and how to put a durable identity on the web.
Treat it as a checklist to work through, not an essay to read once.
\end{minipage}

\vfill
{\color{mutegray} October 2026}\\[0.3cm]
{\footnotesize\color{mutegray} Deadlines, prices, and tool features change.
Verify dates at the source before you plan around them.}
\end{titlepage}

%==============================================================================
\tableofcontents
\thispagestyle{fancy}
\clearpage

%==============================================================================
\section{How to use this guide}

A PhD is won on compounding habits, not on heroic sprints. The first three
months decide how much friction every later month carries: whether you can find
a paper you read in March when you need it in November, whether your experiments
are reproducible, whether your name is findable, whether a deadline surprises you.

Work through this document in order the first time. Section~\ref{sec:90days}
is the short version.

\subsection{The first 90 days, condensed}
\label{sec:90days}

\begin{cbox}{Week 1 --- infrastructure}
\begin{itemize}
  \item Install \tool{Zotero} + the browser connector + the \tool{Better BibTeX}
        plugin (\S\ref{sec:zotero}).
  \item Create a \tool{Google Scholar} profile and set up three alerts
        (\S\ref{sec:scholar}).
  \item Create accounts: \tool{Overleaf}, \tool{OpenReview},
        \tool{GitHub}, \tool{Hugging Face}, \tool{arXiv}, \tool{ORCID}
        (\S\ref{sec:writing}).
  \item Bookmark the AI deadlines countdown and subscribe the relevant
        deadlines to your calendar (\S\ref{sec:deadlines}).
  \item Set up your AI assistant with a project memory file
        (\S\ref{sec:assistant}).
  \item Create one git repository per project, with a \code{notes/} lab
        notebook from day one (\S\ref{sec:infra}).
\end{itemize}
\end{cbox}

\begin{cbox}{Weeks 2--10 --- a focused survey}
Pick a narrow question with your advisor and run a small systematic survey
(\S\ref{sec:survey}). Output: a 60--120 paper annotated bibliography, an
extraction spreadsheet, a taxonomy figure, 10--15 pages of writing, and a
shortlist of three gaps that could become your first paper. This is the single
highest-return thing you can do in your first term.
\end{cbox}

\begin{cbox}{Weeks 10--13 --- identity and first experiments}
\begin{itemize}
  \item Publish \code{username.github.io} with the AcademicPages template
        (\S\ref{sec:website}).
  \item Reproduce one baseline from the survey end to end, on your own
        hardware, with tracked runs (\S\ref{sec:infra}).
  \item Pick a target venue and work backwards from its deadline
        (\S\ref{sec:deadlines}).
\end{itemize}
\end{cbox}

%==============================================================================
\section{The publication landscape}
\label{sec:venues}

ML and NLP are \emph{conference-first} fields. Unlike most of science, the
top-tier results appear at conferences with full peer review and archival
proceedings; journals play a secondary (though rising) role. Your CV will be
read as a list of venue acronyms, so learn them early.

Three things to internalise:

\begin{enumerate}
  \item \textbf{Cycles, not rolling submissions.} Most venues have one or two
        hard deadlines per year. Miss by an hour and you wait months.
  \item \textbf{Reviewing is public-ish.} ICLR and many workshops review on
        OpenReview, where reviews and rebuttals stay online. Read them --- they
        are the best free writing lesson available.
  \item \textbf{A paper has a lifecycle:} preprint $\rightarrow$ submission
        $\rightarrow$ reviews $\rightarrow$ rebuttal $\rightarrow$ decision
        $\rightarrow$ camera-ready $\rightarrow$ talk/poster $\rightarrow$
        citations. Plan for 6--12 months from idea to published.
\end{enumerate}

\subsection{Core machine learning venues}

\begin{table}[h]
\centering
\small
\begin{tabularx}{\textwidth}{L{2.4cm} X L{3.5cm}}
\toprule
\thead{Venue} & \thead{Scope} & \thead{Rhythm (verify yearly)} \\
\midrule
\rowcolor{rowtint}
\textbf{NeurIPS} & Broadest ML venue: theory, deep learning, datasets \&
benchmarks track, many workshops & Abstract $\sim$May, full $\sim$May;
conference December \\
\textbf{ICML} & Core ML, strong on methods and theory &
Deadline $\sim$Jan--Feb; conference July \\
\rowcolor{rowtint}
\textbf{ICLR} & Representation learning, deep learning, LLMs; fully open
review on OpenReview & Deadline $\sim$Sept; conference April--May \\
\textbf{AISTATS} & Statistics $\times$ ML, principled methods &
Deadline $\sim$Oct; conference April--May \\
\rowcolor{rowtint}
\textbf{COLT / UAI} & Learning theory / probabilistic and causal reasoning &
Winter--spring deadlines \\
\textbf{AAAI / IJCAI} & Broad AI, including applications and multi-agent &
AAAI $\sim$Aug; IJCAI $\sim$Jan \\
\rowcolor{rowtint}
\textbf{COLM} & Dedicated to language modelling: data, training, evaluation,
science of LMs & Deadline $\sim$March; conference October \\
\textbf{MLSys / LoG} & Systems for ML / learning on graphs &
Autumn deadlines \\
\bottomrule
\end{tabularx}
\caption*{\small\color{mutegray} Dates are indicative patterns, not promises.
Always confirm on the venue site and the deadline tracker (\S\ref{sec:deadlines}).}
\end{table}

\subsection{Core NLP venues and the ACL Rolling Review}

\begin{table}[h]
\centering
\small
\begin{tabularx}{\textwidth}{L{2.4cm} X L{3.5cm}}
\toprule
\thead{Venue} & \thead{Scope} & \thead{Notes} \\
\midrule
\rowcolor{rowtint}
\textbf{ACL} & The flagship: all of computational linguistics and NLP &
Summer conference; via ARR commitment \\
\textbf{EMNLP} & Empirical methods --- the other flagship &
Late-year conference; via ARR \\
\rowcolor{rowtint}
\textbf{NAACL / EACL / AACL} & Regional ACL conferences (Americas, Europe,
Asia-Pacific); same standard & Alternate by year \\
\textbf{COLING / LREC} & Computational linguistics / language resources and
evaluation & Historically biennial; check the current cycle \\
\rowcolor{rowtint}
\textbf{CoNLL / *SEM / SemEval} & Learning and language, semantics, shared
tasks & Good first-paper targets \\
\textbf{Interspeech / ICASSP} & Speech and audio & Speech-specific cycles \\
\rowcolor{rowtint}
\textbf{TACL} & Journal of the ACL --- rolling submission, presented at *ACL
conferences & Slower but prestigious \\
\textbf{Findings} & Companion archival track of *ACL venues for solid papers
below the main-conference bar & A real, citable publication \\
\bottomrule
\end{tabularx}
\end{table}

\begin{cbox}{ARR: the mechanic that confuses every first-year NLP student}
The ACL family separates \emph{reviewing} from \emph{publishing}. You submit to
\textbf{ACL Rolling Review} (ARR) in one of its cycles --- every two to three
months in 2026, not monthly; check the current calendar. You
receive reviews plus a meta-review. Then, in a separate step, you
\emph{commit} that reviewed paper to a specific venue (ACL, EMNLP, NAACL,
\ldots), whose chairs decide accept / Findings / reject based on those reviews.

Consequences worth planning around:
\begin{itemize}
  \item There are two deadlines to track: the ARR \emph{submission} deadline
        and the venue \emph{commitment} deadline.
  \item Reviews are reusable --- a rejected commitment can be revised and
        committed elsewhere, or resubmitted to ARR with a revision note.
  \item Mandatory sections: a \textbf{Limitations} section (does not count
        against the page limit) and the \textbf{Responsible NLP checklist}.
        Missing limitations is a desk-reject risk.
\end{itemize}
\end{cbox}

\subsection{Adjacent venues and journals}

\begin{itemize}
  \item \textbf{Data mining / web / IR:} KDD, WSDM, CIKM, SIGIR, The Web
        Conference, ECML-PKDD. Natural homes for applied and graph work.
  \item \textbf{Vision:} CVPR, ICCV, ECCV --- relevant for multimodal work.
  \item \textbf{Journals:} JMLR (open access, strong), \textbf{TMLR}
        (rolling, fast, judged on correctness + interest rather than
        novelty hype --- an excellent fit for careful work that conference
        reviewers undervalue), TACL, \emph{Computational Linguistics}, TPAMI.
  \item \textbf{Workshops:} co-located, non-archival or lightly archival,
        4--8 pages, fast turnaround, real feedback and real audiences
        (BlackboxNLP, RepL4NLP, WMT, and dozens more each year). The honest
        advice for a first-year student: \emph{aim your first submission at a
        good workshop}. You learn the full pipeline --- writing, reviewing,
        rebuttal, poster, conversations --- at a fraction of the risk.
\end{itemize}

\subsection{How to judge a venue without asking around}

\begin{itemize}
  \item \textbf{Google Scholar Metrics} --- per-venue h5-index
        (\url{https://scholar.google.com/citations?view_op=top_venues}).
  \item \textbf{CORE ranking} (\url{https://portal.core.edu.au/conf-ranks/})
        --- A*/A/B/C tiers; used for evaluation in many countries.
  \item \textbf{CSRankings} (\url{https://csrankings.org}) --- shows which
        venues count as top-tier per subarea, which is the useful signal.
  \item \textbf{Where does the work you admire get published?} The most
        reliable test. Open the last three papers you loved and look at the
        venue line.
\end{itemize}

\begin{wbox}{Predatory venues}
If a conference emails you an invitation, promises acceptance in two weeks,
charges a large fee, and spans ``AI, Big Data, Blockchain and IoT'' --- it is a
scam. Publishing there actively damages your CV. Cross-check every unfamiliar
venue against the sources above before submitting.
\end{wbox}

%==============================================================================
\section{Deadlines: never be surprised again}
\label{sec:deadlines}

\begin{cbox}{The one page to bookmark}
Conference deadlines are countdowns, not dates: they are stated in
\textbf{AoE} (Anywhere on Earth $=$ UTC$-$12), so an ``11:59pm'' deadline is
later than you think --- and a venue has far more dates than the one everybody
remembers. Two boards cover this, and they complement each other.

\begin{itemize}
  \item \textbf{Papers with Code --- AI Deadlines} (start here):\\
        \url{https://paperswithcode.co/ai-deadlines}\\
        Filter by subfield (Machine Learning, NLP,
        Computer Vision, Data Mining, Robotics, IR, Speech, RL, \ldots), or
        restrict to \emph{top tier only}. Crucially it tracks \emph{every}
        milestone, not just the submission date: abstract registration,
        reviews released, author-feedback and rebuttal windows, commitment
        dates, notifications, camera-ready --- each with a live countdown,
        the conference dates, and the location. The same site carries trending
        papers, a following feed, and an API and MCP endpoint you can wire
        into your assistant (\S\ref{sec:assistant}).
  \item \textbf{AI Deadlines}, pre-filtered for your subfields:\\
        \url{http://aideadlines.org/?sub=ML,CV,NLP}\\
        The \code{sub=} parameter carries the filter, so the link itself is
        your personalised board --- bookmark it and edit the tags
        (\code{ML}, \code{CV}, \code{NLP}, \code{SP}, \code{DM}, \code{RO},
        \ldots) as your interests change. A mirror runs as a Hugging Face
        Space (\url{https://huggingface.co/spaces/huggingface/ai-deadlines}),
        and the data is an open repository of YAML files
        (\url{https://github.com/huggingface/ai-deadlines}) --- so a missing
        deadline is a deadline you can add by pull request, a trivially easy
        first open-source contribution.
\end{itemize}
Cross-check both against the official venue page before you plan around a
date. Aggregators lag; call-for-papers pages do not.
\end{cbox}

A workflow that actually holds:

\begin{enumerate}
  \item Filter the board to your subfields (NLP, ML, maybe IR/CV) and list the
        4--6 venues you could plausibly target in the next 18 months.
  \item Put each into your calendar \emph{three times}: the deadline itself,
        minus 7 days (``results frozen, writing only''), and minus 21 days
        (``experiments must be running'').
  \item Add the secondary dates people forget: abstract deadlines (NeurIPS),
        ARR commitment dates, supplementary-material deadlines, rebuttal
        windows (often a hard 5--7 day window --- do not travel then),
        camera-ready dates, and early-registration dates.
  \item Backward-plan on paper: camera-ready $\leftarrow$ decision $\leftarrow$
        rebuttal $\leftarrow$ submission $\leftarrow$ freeze $\leftarrow$ final
        experiments $\leftarrow$ ablations $\leftarrow$ main result.
        Ablations always take longer than the main result.
  \item Alternative trackers to cross-check: \url{https://ccfddl.com} and the
        official venue pages. Trust the venue page over any aggregator.
\end{enumerate}

\begin{wbox}{Deadline-driven research is a trap, in moderation}
Deadlines are useful forcing functions and terrible research strategies. Keep
one long-horizon question alive alongside the deadline work --- it is where
your thesis contribution will come from.
\end{wbox}

%==============================================================================
\section{Writing and sharing: LaTeX and Overleaf}
\label{sec:writing}

Every paper you write in this field will be LaTeX. Learn it properly once ---
a weekend --- and it will never cost you again. What differs is \emph{where}
you edit and \emph{how} you collaborate. (Typst is pleasant and growing, but
no venue template accepts it, so papers stay in LaTeX.)

\subsection{Overleaf --- the default collaborative editor}

\url{https://www.overleaf.com}

\begin{itemize}
  \item Browser-based, nothing to install, compiles server-side. Co-authors
        type in the same document in real time, which is what you want in the
        48 hours before a deadline.
  \item Template gallery contains the official style files for most venues
        --- search ``NeurIPS'', ``ACL'', ``ICML'' before you start a paper and
        begin from the official template, never from last year's copy.
  \item Rich comments and a chat panel; useful for advisor feedback.
  \item Overleaf now ships its own \textbf{AI Assistant} inside the editor:
        LaTeX-aware, and metered by plan --- a handful of uses a day on the
        free tier, more on paid ones, unmetered at the top. The cautions
        below apply to it word for word.
  \item Many universities have institutional licences --- check before paying.
        The free tier is limited on collaborator count, compile time, full
        version history, Git/Dropbox sync, and track changes.
  \item Share with a link (view or edit), which is how most advisors want to
        receive a draft.
\end{itemize}

\subsection{Prism --- closing on 1 January 2027}

\begin{wbox}{Do not start anything in Prism}
OpenAI's AI-native LaTeX workspace (\url{https://prism.openai.com}, launched
January 2026) \textbf{shuts down on 1 January 2027}. It was free and genuinely
good --- a frontier model that saw the whole manuscript rather than isolated
snippets --- and it lasted under a year.

\begin{itemize}
  \item \textbf{If you have projects there, export them now.} The site carries
        an export link; do not leave it to the last week.
  \item \textbf{What exports do \emph{not} carry:} comments, version history,
        collaborators and linked integrations. If a co-author's comments
        matter, copy them out by hand before the date.
  \item \textbf{Land somewhere boring:} Overleaf for collaboration, or
        \code{git} + \code{latexmk} locally. Both will outlive your PhD.
\end{itemize}
\end{wbox}

\begin{cbox}{The transferable lesson}
A free cloud tool can disappear with a quarter's notice, and your thesis
cannot. Keep every manuscript in \textbf{plain \code{.tex} files under git}
that you can compile offline, and treat any editor --- Overleaf included ---
as a convenient window onto that repository rather than as the place your work
lives.
\end{cbox}

\begin{wbox}{Two cautions before you move your thesis into any AI editor}
\textbf{Policy:} most venues now have an explicit policy on AI-assisted
writing (disclosure requirements, and a hard ban on using LLMs to write
\emph{reviews}). You remain fully responsible for every sentence, citation,
and number. Read your venue's policy once, now.\\[4pt]
\textbf{Confidentiality:} do not upload data, text, or results covered by an
NDA, an industry collaboration, or someone else's unpublished work into a
cloud tool without checking. This is a real way for PhD students to get into
real trouble.
\end{wbox}

\subsection{The local workflow you should also have}

Cloud editors are for collaboration; a local copy is for control. Keep both.

\begin{itemize}
  \item \code{git} + \code{latexmk -pdf -pvc} + VS Code with the LaTeX
        Workshop extension (or your editor's equivalent).
  \item \textbf{One sentence per line} in your \code{.tex} files. This single
        convention turns git diffs from useless to invaluable.
  \item Split by section (\code{\textbackslash input\{sections/method\}}) so
        two people can work without conflicts.
  \item Keep figures as generated artefacts: a script in the repo produces
        every PDF figure from logged results. No hand-edited plots.
  \item \textbf{Textual vs.\ parenthetical citations.} With natbib (the ACL
        and NeurIPS templates): \code{\textbackslash citet\{\}} when the
        authors are part of the sentence (``Devlin et al.\ (2019) show''),
        \code{\textbackslash citep\{\}} when the citation is an aside
        (``(Devlin et al., 2019)''). With biblatex: \code{\textbackslash
        textcite} and \code{\textbackslash parencite}. Mixing them up is the
        most common formatting error in student drafts.
  \item \textbf{\code{latexdiff}} (on CTAN) marks up every change between two
        versions of the source: ideal for a camera-ready, a journal revision,
        or showing your advisor what moved since last week.
  \item For arXiv submission, \code{arxiv-latex-cleaner}
        (\,\url{https://github.com/google-research/arxiv-latex-cleaner}\,)
        strips comments, removes unused files, and shrinks images --- run it
        every time, not least because \emph{reviewers have found deanonymising
        information in LaTeX comments}.
\end{itemize}

\subsection{Preprints, anonymity, and sharing}

\begin{itemize}
  \item \textbf{arXiv} (\url{https://arxiv.org}) --- cs.CL for NLP, cs.LG for
        ML, cs.AI, stat.ML. Your first submission needs endorsement; your
        advisor or a co-author handles it. Preprints are the norm and
        accelerate citations.
  \item \textbf{Anonymity periods:} ACL-family venues enforce an anonymity
        window around the review cycle --- posting or publicising a new
        preprint during it can invalidate your submission. ML venues
        (NeurIPS/ICML/ICLR) are double-blind but generally permit preprints.
        \emph{Check the current policy for your target venue before you post.}
  \item For double-blind submissions, anonymise your code too:
        \url{https://anonymous.4open.science}. Also strip author names from
        PDF metadata and acknowledgements.
  \item \textbf{Reproducibility artefacts:} release code on GitHub with a
        licence, checkpoints and datasets on Hugging Face, and an
        exact environment spec. Papers with runnable code get cited more ---
        and your future self is the main beneficiary.
  \item \textbf{Make the code citable.} A \code{CITATION.cff} file in the
        repo (\url{https://citation-file-format.github.io}) gives GitHub a
        ``Cite this repository'' button, and the Zenodo integration
        (\url{https://zenodo.org}) mints a DOI for every tagged release, so
        the code can be cited independently of the paper.
  \item Get an \textbf{ORCID} (\url{https://orcid.org}) in week one and put it
        on everything; it disambiguates you from everyone who shares your name.
\end{itemize}

%==============================================================================
\section{Your reference library: Zotero}
\label{sec:zotero}

You will read a few thousand papers during a PhD. Without a system you will
re-find the same paper five times and mis-cite it once in print. Set this up
before you read paper number ten.

\url{https://www.zotero.org} --- free, open source, no lock-in.

\subsection{The setup, in order}

\begin{enumerate}
  \item Install \textbf{Zotero} (desktop) and the \textbf{Zotero Connector}
        browser extension. The connector is the whole point: one click on an
        arXiv, ACL Anthology, OpenReview, or Scholar page saves the metadata
        \emph{and} the PDF.
  \item Install \textbf{Better BibTeX} (\url{https://retorque.re/zotero-better-bibtex/}).
        This is non-optional for LaTeX users. It gives you:
        \begin{itemize}
          \item stable, readable citation keys (\code{devlin2019bert}) that you
                \textbf{pin} so they never change under you;
          \item \textbf{auto-export}: point it at \code{paper/refs.bib} in your
                git repo and it rewrites that file whenever your library
                changes. Your bibliography is then always current and you never
                hand-edit BibTeX again.
        \end{itemize}
  \item Turn on sync (free account; 300\,MB of free file storage, paid tiers or
        your own WebDAV/cloud folder beyond that). Attachments can also be
        stored in a linked folder you back up yourself.
  \item Create a \textbf{group library} with your lab or co-authors --- shared
        reading lists, shared annotations, one source of truth per project.
  \item Organise with \textbf{collections} per project and \textbf{tags} per
        theme (\code{method:retrieval}, \code{status:to-read},
        \code{baseline}). Saved searches then do the filing for you.
  \item Read inside Zotero: its built-in PDF reader keeps highlights and
        notes attached to the item, so your annotations are searchable years
        later.
\end{enumerate}

\subsection{Getting a reference into your \code{.bib}}

\begin{cbox}{The easiest route, and it needs no setup}
\begin{enumerate}
  \item In Google Scholar: \emph{Settings} $\rightarrow$ \emph{Bibliography
        manager} $\rightarrow$ tick \emph{Show links to import citations into
        BibTeX}. One time only.
  \item A \textbf{Cite} link now appears under every search result. Click it,
        then \textbf{BibTeX}, and copy the entry.
  \item Paste it into your \code{refs.bib} --- or, better, into Zotero with
        \emph{File $\rightarrow$ Import from clipboard}, so your library and
        your bibliography stay in step.
\end{enumerate}
Scholar also exports in bulk: save papers to \emph{My library} with the star
as you find them, then select them and \emph{Export} $\rightarrow$
\emph{BibTeX}.
\end{cbox}

\begin{wbox}{Whatever Scholar gives you, check it}
Expect \code{arXiv preprint arXiv:1810.04805} for a paper that was published
at NAACL, missing pages and editors, venue names that differ between entries,
and mangled accents in author names. Check \code{booktitle}/\code{journal},
year and pages before you trust an entry --- a bibliography full of ``arXiv
preprint'' is the clearest signal that nobody looked.
\end{wbox}

Cleaner sources, each one click away:

\begin{itemize}
  \item \textbf{ACL Anthology} --- every paper page has a BibTeX box, and
        every volume offers one \code{.bib} for the entire proceedings.
  \item \textbf{DBLP} --- per-paper BibTeX, plus a one-click export of
        everything an author has published.
  \item \textbf{arXiv} --- ``Export BibTeX citation'' on the abstract page
        (fine for genuinely unpublished work).
  \item \textbf{Semantic Scholar}, the ACM Digital Library, IEEE Xplore and
        Springer all have a cite or export button.
  \item \textbf{Any DOI, straight from the terminal:}\\
        \code{curl -LH "Accept: application/x-bibtex" https://doi.org/<DOI>}\\
        returns the publisher's own entry --- correct venue, pages and year.
        Ideal in a script, or when a site buries its export button.
\end{itemize}

Past twenty or so references, stop copy-pasting entirely: the Zotero connector
plus Better BibTeX auto-export does all of the above for you and keeps your
keys stable.

\subsection{Habits that keep a library trustworthy}

\begin{itemize}
  \item \textbf{Cite the published version, not the preprint}, once one
        exists. Reviewers notice; it is a cheap signal of care. Keep both
        entries if you must, but cite the proceedings.
  \item \textbf{Get clean BibTeX from the right source.} For CS, the ACL
        Anthology (\url{https://aclanthology.org}) and
        \textbf{DBLP} (\url{https://dblp.org}) produce correct proceedings
        entries; Google Scholar's BibTeX is often sloppy about venue names,
        pages, and editors. Zotero + Add Item by Identifier (DOI / arXiv ID /
        ISBN) is the fastest clean path.
  \item Run \emph{Duplicate Items} monthly. Merge, do not delete.
  \item \textbf{Retraction alerts are already on.} Zotero checks your library
        against Retraction Watch and marks retracted items with a red banner;
        when you see one, find out why before you cite it.
  \item \textbf{Tidy the file, not the entries.} \code{bibtex-tidy}
        (\url{https://github.com/FlamingTempura/bibtex-tidy}; there is also a
        web version) formats, sorts, deduplicates and strips junk fields such
        as \code{abstract} and \code{file}. One command before every
        submission --- or a Makefile target.
  \item \textbf{Audit before you submit.} A reference that does not resolve
        is a typo or a fabrication, and reviewers now look for both. Wrap the
        \code{curl} line above in a loop over every DOI in \code{refs.bib}
        (and the arXiv API for \code{eprint} fields), flag anything that
        returns no title, and open whatever is flagged. Ten lines of shell or
        Python; ask your assistant to write them and keep the script in the
        repo.
  \item Add one-sentence notes on save: ``what it does / why I care''. In
        month 14 this is the difference between a library and a junkyard.
  \item Subscribe to feeds inside Zotero (arXiv category RSS, journal feeds) so
        new work arrives where you already work.
\end{itemize}

\begin{wbox}{Three BibTeX gotchas that survive every tool}
\textbf{Capitals vanish.} Most styles lowercase titles, so protect what must
stay upper-case: \code{title = \{\{BERT\}: Pre-training of ...\}}, and the
same for \code{\{GPT-4\}}, \code{\{Greek\}}, \code{\{LLM\}}.\\[3pt]
\textbf{Author format.} \code{author = \{Devlin, Jacob and Chang, Ming-Wei\}}
--- \code{Last, First}, joined by \code{and}; never a comma-separated list,
never an ``et al.'' typed by hand.\\[3pt]
\textbf{arXiv-only work} gets the arXiv fields, not an invented journal:
\code{@misc\{..., eprint=\{2505.24403\}, archivePrefix=\{arXiv\},
primaryClass=\{cs.LG\}\}}. And pick \emph{one} source per venue family:
DBLP's long booktitles next to the Anthology's short ones look careless.
\end{wbox}

\begin{cbox}{Zotero $\times$ Overleaf}
Overleaf's paid tiers can link a Zotero library directly. On a free plan, the
Better BibTeX auto-export route is better anyway: your \code{.bib} lives in
git, so it is versioned, diffable, and identical for every co-author. Commit it.
\end{cbox}

%==============================================================================
\section{Your first real move: a focused survey}
\label{sec:survey}

Before you can contribute, you need to know what exists --- and the fastest
honest way to learn a field is to survey it systematically. Not reading
``around'' the topic: producing a structured artefact.

\begin{cbox}{Why this is the right first project}
\begin{itemize}
  \item It gives you the \textbf{vocabulary} to talk to your advisor as a peer
        within weeks.
  \item It produces a \textbf{taxonomy} --- and a taxonomy is the thing that
        reveals gaps. Your first paper usually lives in a box you drew that
        has nothing in it.
  \item It builds your \textbf{bibliography}, which you will reuse in every
        paper and in your thesis introduction.
  \item It identifies the \textbf{baselines} you will have to beat, and the
        datasets and metrics everyone uses, before you waste a month on a
        non-standard setup.
  \item It is \textbf{publishable}: as a survey, a position paper, or a
        workshop paper --- and it is the natural core of your thesis'
        related-work chapter.
\end{itemize}
\end{cbox}

\subsection{A protocol you can run in 8--12 weeks}

\begin{description}
  \item[1. Write the question down.] One sentence, narrow enough to be
        answerable: not ``LLM reasoning'' but ``methods for verifying
        intermediate reasoning steps in LLMs, 2022--2026, and how they are
        evaluated''. Agree it with your advisor in writing.
  \item[2. Seed set.] Ask your advisor for 5--10 anchor papers. Read them
        properly. Note their shared vocabulary --- those terms become your
        search queries.
  \item[3. Define inclusion / exclusion criteria.] E.g.\ include: peer-reviewed
        or arXiv, 2022+, English, proposes or evaluates a verification method.
        Exclude: pure position pieces, no empirical evaluation. Write the
        criteria down \emph{before} screening, or you will rationalise.
  \item[4. Search systematically.] Run the same queries across Google Scholar,
        Semantic Scholar, the ACL Anthology, OpenReview, and DBLP
        (\S\ref{sec:scholar}). \textbf{Log every query, the date, and the
        number of hits} in a spreadsheet tab. This is what makes a survey
        reproducible rather than anecdotal --- and it is the discipline the
        PRISMA standard formalises in other fields.
  \item[5. Snowball in both directions.] \emph{Backward:} mine the reference
        lists of your seed papers. \emph{Forward:} for each seed, open
        ``Cited by'' and search \emph{within} the citing articles. Forward
        snowballing is how you find the last six months, which no survey
        covers yet.
  \item[6. Screen in two passes.] Title + abstract first (fast, ruthless,
        keep a ``maybe'' pile), then full text for survivors. Expect to screen
        400--600 to keep 60--120. Past a few hundred, \textbf{ASReview}
        (\url{https://asreview.nl}; free, open source) learns from your first
        decisions and orders the rest by likely relevance; Rayyan is the
        lighter, hosted alternative. Neither replaces the criteria you wrote
        in step 3.
  \item[7. Extract into a table, not prose.] One row per paper, fixed columns
        (Table~\ref{tab:extract}). Fill it as you read. The table \emph{is}
        your survey; the prose comes from reading down its columns.
  \item[8. Build the taxonomy.] Cluster rows until the groups stop changing.
        Draw it as one figure. If you cannot draw it, you do not understand
        the field yet --- which is useful information.
  \item[9. Write 10--15 pages as you go,] not at the end. Each section
        summarises one branch of the taxonomy and ends with ``what is missing
        here''.
  \item[10. Finish with three candidate gaps,] each stated as a testable
        question plus the experiment that would answer it. Take those to your
        advisor. That meeting is where your first paper starts.
\end{description}

\begin{table}[h]
\centering
\small
\begin{tabularx}{\textwidth}{L{3.1cm} X}
\toprule
\thead{Column} & \thead{What goes in it} \\
\midrule
\rowcolor{rowtint}
Key & Pinned Better BibTeX key (\code{author2026short}) \\
Venue / year & Published venue, not ``arXiv'' if it has one \\
\rowcolor{rowtint}
Problem & The task, in your words, in one line \\
Method family & Your taxonomy label (fill this last, revise often) \\
\rowcolor{rowtint}
Data & Datasets and sizes used \\
Metrics & How success is measured; is it comparable across papers? \\
\rowcolor{rowtint}
Headline result & The number they lead with, plus the baseline it beats \\
Compute & GPUs, hours, model sizes --- can \emph{you} run this? \\
\rowcolor{rowtint}
Code & Link + does it reproduce the paper's numbers? (\S\ref{sec:signals}) \\
Limitations & Stated and unstated \\
\rowcolor{rowtint}
Relevance & Your 1--5 score and a one-line ``why I care'' \\
\bottomrule
\end{tabularx}
\caption{Extraction schema. A spreadsheet with these columns, 100 rows deep, is
worth more than 100 pages of highlighted PDFs.}
\label{tab:extract}
\end{table}

\begin{cbox}{Read efficiently: the three-pass method}
Keshav's \emph{How to Read a Paper} is two pages and will save you months.
Pass 1 (5~min): title, abstract, intro, section headings, conclusions,
references --- decide whether to continue. Pass 2 (1~h): figures, tables,
method at a level where you could summarise it. Pass 3 (several hours, rare):
re-derive the work, assumption by assumption. Give pass~3 to about 20 papers;
pass~1 to hundreds.
\end{cbox}

\begin{wbox}{Scope discipline}
A survey expands without limit if you let it. Timebox it, freeze the search
date (``literature up to 15 November 2026''), and ship. You can always write
version two --- and a frozen, documented cut-off is exactly what a reviewer
wants to see.
\end{wbox}

%==============================================================================
\section{Searching the literature properly}
\label{sec:scholar}

Most people use Google Scholar as a search box. It is closer to a research
instrument, and the difference shows up in what you find.

\subsection{Google Scholar: the operators and features that matter}

\url{https://scholar.google.com}

\begin{table}[h]
\centering
\small
\begin{tabularx}{\textwidth}{L{5.0cm} X}
\toprule
\thead{Query / feature} & \thead{What it does} \\
\midrule
\rowcolor{rowtint}
\code{"contrastive decoding"} & Exact phrase. Always quote multi-word technical
terms, or Scholar matches them loosely. \\
\code{retrieval OR retrieve} & \code{OR} must be capitalised. Essential,
because fields rename themselves constantly. \\
\rowcolor{rowtint}
\code{-survey -review} & Excludes a term. Useful once surveys swamp your
results. \\
\code{author:"K Skianis"} & Author search; try initial + surname, and
surname variants. \\
\rowcolor{rowtint}
\code{source:"Empirical Methods"} & Restricts to a publication whose name
contains the phrase --- the way to search within a venue. \\
\code{intitle:hallucination} & Term must appear in the title. Brutal and
effective precision filter; \code{allintitle:} for all terms. \\
\rowcolor{rowtint}
\textbf{Custom range} (left panel) & Year filtering. Combine with
\textbf{Sort by date} to see this month's work rather than the most-cited work
of 2019. \\
\textbf{Cited by $N$} & Forward citations --- who built on this paper. \\
\rowcolor{rowtint}
\textbf{Search within citing articles} & The most under-used feature in
academia: search \emph{inside} the set of papers citing a key work. ``Who cited
BERT \emph{and} talks about Greek?'' is one query. \\
\textbf{Related articles} & Scholar's similarity graph. Good for escaping your
own vocabulary. \\
\rowcolor{rowtint}
\textbf{All $N$ versions} & Finds the free PDF and the published version of the
same work. \\
\textbf{Cite $\rightarrow$ BibTeX} & Enable it once: Settings $\rightarrow$
Bibliography manager $\rightarrow$ Show links to import citations into BibTeX.
(Still prefer DBLP/ACL Anthology for final entries.) \\
\bottomrule
\end{tabularx}
\end{table}

\subsubsection*{Scholar Labs --- ask the question, not the keywords}

\url{https://scholar.google.com/scholar_labs/search} --- a search mode inside
Scholar, launched in November 2025 and open to all signed-in users, in any
language, since June 2026. Instead of a keyword string you type the actual
research question --- ``which methods verify intermediate reasoning steps in
LLMs, and how are they evaluated?'' --- and it breaks the question into its
topics and relationships, searches each, and returns papers that answer the
whole question, with a sentence on \emph{how} each one does and follow-up
questions to narrow in. Use it in the first week of a survey, when you do not
yet know the field's vocabulary; log its queries like any other
(\S\ref{sec:survey}); and read the papers --- the explanation is a lead, not a
summary you can cite.

\subsubsection*{Alerts --- set these up today}

\begin{itemize}
  \item \textbf{Query alerts} for your 2--3 core topics (use quoted phrases, or
        you will drown).
  \item \textbf{Citation alerts} on the 5 anchor papers of your area --- new
        work almost always cites them, so this is a near-complete feed of your
        subfield.
  \item \textbf{An alert on your own profile}, so you learn who is citing you.
  \item Set a weekly slot (30 minutes, same time) to triage the alert emails
        into Zotero. Alerts you never read are worse than no alerts.
\end{itemize}

\subsubsection*{Your Scholar profile}

Create it in week one, even with zero papers. Use your institutional email so
it gets verified, add a photo, list your interests, keep the publication list
clean (merge duplicates, delete the garbage Scholar attributes to you). Decide
whether to make it public --- for a researcher, public is almost always right.
It is how programme chairs, reviewers, and future employers check who you are.

\begin{wbox}{On h-index and citation counts}
They are age-dependent, field-dependent, and gameable. A first-year student
with $h=0$ is exactly on schedule. Use these numbers as \emph{navigation aids
for finding influential work}, never as a measure of your worth or anyone
else's.
\end{wbox}

\subsection{The rest of the toolkit}

\begin{description}
  \item[Semantic Scholar] (\url{https://www.semanticscholar.org}) --- better
        structured than Scholar: TLDR summaries, ``highly influential
        citations'', citation \emph{intent} (background vs.\ method vs.\
        result), and a genuinely good free API for building your own tooling.
  \item[OpenAlex] (\url{https://openalex.org}) --- fully open bibliographic
        graph with a free API. The right tool if you want to compute anything
        over the literature (growth curves, co-authorship, venue trends).
  \item[ACL Anthology] (\url{https://aclanthology.org}) --- every *ACL paper
        ever, free, full-text searchable, with correct BibTeX. For NLP it
        beats Scholar for precision.
  \item[OpenReview] (\url{https://openreview.net}) --- submissions,
        \emph{reviews}, rebuttals, and decisions for ICLR, NeurIPS and many
        workshops. Read the reviews of papers adjacent to yours before you
        write your own. Nothing else teaches you the standard this quickly.
  \item[DBLP] (\url{https://dblp.org}) --- authoritative CS metadata. Use it
        to enumerate everything a person or venue published, and for clean
        BibTeX.
  \item[Connected Papers] (\url{https://www.connectedpapers.com}),
        \textbf{ResearchRabbit}, \textbf{Litmaps} --- visual citation
        neighbourhoods from one seed paper. Excellent for the first week of a
        survey, when you do not yet know the right words.
  \item[Scite] (\url{https://scite.ai}) --- tags each citation as
        supporting, mentioning or contrasting the cited claim, which is the
        question citation \emph{counts} cannot answer. Paid beyond a trial;
        worth it for a survey.
  \item[Unpaywall] (\url{https://unpaywall.org/products/extension}) --- a
        browser extension that finds the legal open-access copy of a paywalled
        paper. Install it the day your institutional access fails you.
  \item[Hugging Face] (\url{https://huggingface.co/papers}) --- daily papers,
        plus models and datasets with download counts: a usage signal that
        citation counts miss entirely.
  \item[alphaXiv / arXiv listings] --- discussion layers and raw daily feeds
        for cs.CL and cs.LG. Skim titles daily, 10 minutes, no guilt about
        what you skip.
\end{description}

\begin{cbox}{Papers with Code, then and now}
The original \code{paperswithcode.com} was discontinued in 2025; its domain
now redirects to Hugging Face's own Papers hub, so old links still work but
land there. A separate successor under the Papers with Code name runs at
\url{https://paperswithcode.co}: trending papers, a following feed for tasks,
authors and organisations, a paper archive, the deadline board of
\S\ref{sec:deadlines}, and an API plus MCP endpoint. The site does not say who
operates it, so treat it as a convenient mirror and cite original sources.
\end{cbox}

%==============================================================================
\section{Reading the signals: citations, stars, and attention}
\label{sec:signals}

Part of research taste is judging, quickly, whether a line of work is gaining
momentum or quietly dying. Two cheap signals do most of the work --- as long as
you know what each actually measures.

\subsection{Citations: measure velocity, not volume}

\begin{itemize}
  \item \textbf{Normalise by age.} A 2019 paper with 400 citations and a 2026
        paper with 60 are not comparable. Compute citations per year since
        publication; for recent work, citations in the \emph{last} 12 months
        is the sharper signal.
  \item \textbf{Look at the second derivative.} A paper whose citations are
        accelerating marks a live area. One with a flat tail marks a closed
        one. Scholar's per-year bar chart on a paper's page shows this at a
        glance.
  \item \textbf{Read \emph{who} cites, not how many.} Ten citations from the
        three strongest groups in the area beat 200 from unrelated surveys.
        Semantic Scholar's citation-intent labels help: being cited as a
        \emph{method} others use is worth far more than a background mention.
  \item \textbf{Watch for inflation:} large author lists citing each other,
        survey-driven citation piles, and benchmark papers that collect
        citations for existing rather than for being good ideas.
  \item \textbf{Venue trend as a field signal.} Count submissions or accepted
        papers mentioning a term per year (OpenReview and the ACL Anthology
        make this scriptable). A term going from 5 to 300 papers in two years
        tells you both that it is hot and that you are probably late.
\end{itemize}

\subsection{GitHub: stars, and the signals that matter more}

Stars are the field's informal applause. They are a real signal of attention
--- and a poor signal of quality.

\begin{table}[h]
\centering
\small
\begin{tabularx}{\textwidth}{L{4.6cm} X}
\toprule
\thead{What to check} & \thead{What it tells you} \\
\midrule
\rowcolor{rowtint}
\textbf{Star history} (\url{https://star-history.com}) & The \emph{shape}
matters: a vertical launch-day spike then flat means one viral post; a steady
upward slope means real adoption. \\
\textbf{Last commit / commit frequency} & 3k stars and nothing pushed in
14 months is an abandoned repo. Check the default branch, not the releases
page. \\
\rowcolor{rowtint}
\textbf{Contributors} & One author vs.\ thirty. Thirty means a community will
keep it alive --- and that bugs get found. \\
\textbf{Open vs.\ closed issues, and who answers} & Authors replying within
days is the strongest health signal there is. \\
\rowcolor{rowtint}
\textbf{Issues searched for ``reproduce''} & Do other people get the paper's
numbers? Search the issue tracker for \code{reproduce}, \code{results},
\code{mismatch} before you trust a baseline. \emph{Do this every time.} \\
\textbf{Forks, ``Used by'', dependents} & Forks are often bookmarks; genuine
downstream dependents are the stronger signal. \\
\rowcolor{rowtint}
\textbf{Releases, tests, CI, pinned deps} & Signals that running it will cost
you a day, not a fortnight. A lockfile is a gift. \\
\textbf{Checkpoints provided} & Can you evaluate without retraining? Often the
difference between a usable and a theoretical baseline. \\
\rowcolor{rowtint}
\textbf{LICENSE} & Research-only, non-commercial, or custom model licences
restrict what you may publish or build on. Check \emph{before} you invest. \\
\bottomrule
\end{tabularx}
\end{table}

\begin{itemize}
  \item Useful GitHub search syntax:
        \code{transformers stars:>500 pushed:>2026-01-01 language:python}, and
        \code{path:} / \code{in:readme} to find code inside repositories.
  \item \textbf{Watch $\rightarrow$ Custom $\rightarrow$ Releases} on the 5--10
        repositories you depend on. Far better than re-checking manually.
  \item Cross-check with other attention signals: Hugging Face model downloads
        and likes, PyPI download counts, and whether the method appears in
        other people's baselines --- the last one is the ground truth.
\end{itemize}

\begin{wbox}{The trap, stated plainly}
Stars measure how well a project was marketed. Citations measure how well a
paper was positioned. Neither measures whether the idea is correct. The only
test that counts is running it yourself on a task you understand --- which is
why reproducing one baseline (\S\ref{sec:90days}) is in your first 90 days.
\end{wbox}

%==============================================================================
\section{Set up an AI assistant --- and the rules for using it}
\label{sec:assistant}

An AI coding and reading assistant is now part of a competent research setup,
in the same way a reference manager is. Used well, it removes weeks of
mechanical work per year. Used badly, it fabricates citations into your
literature review and hollows out the understanding your thesis is supposed to
demonstrate.

\subsection{The setup}

\begin{enumerate}
  \item \textbf{A terminal/IDE coding agent} in each project repository
        (Claude Code, Codex CLI, Cursor, or GitHub Copilot --- any of them).
        This is where the leverage
        is: it reads your code, runs your tests, and edits files.
  \item \textbf{A project memory file} (e.g.\ \code{CLAUDE.md} at the repo
        root) stating your conventions: Python version and env manager, how to
        launch training, where data lives, plotting style, ``never commit
        checkpoints'', ``one sentence per line in .tex''. Written once, it
        stops you re-explaining your project every session.
  \item \textbf{Connect your sources.} Via connectors/MCP servers, give the
        assistant access to your filesystem, your Zotero library, and
        literature APIs (arXiv, Semantic Scholar). An assistant that can read
        the actual PDFs in your library is a different tool from one that
        guesses.
  \item \textbf{Saved prompts for repeated jobs} --- paper triage, extraction
        into your survey schema, ``review this draft as a hostile ARR
        reviewer''. Keep them in the repo so they improve over time.
  \item \textbf{A lab notebook it can read and write:} \code{notes/} with one
        dated Markdown file per week. Decisions, failed runs and why,
        hyperparameters, links. This is also what your assistant needs in
        order to be useful in month six.
\end{enumerate}

\subsection{What to use it for}

\begin{multicols}{2}
\begin{itemize}
  \item Triaging abstracts against your inclusion criteria.
  \item Pulling tables/numbers out of PDFs into your extraction sheet.
  \item Explaining unfamiliar notation; walking a derivation step by step.
  \item Getting an unfamiliar baseline repo to run; writing eval harnesses.
  \item Data pipelines, plotting code, Slurm scripts, boilerplate.
  \item Refactoring research code before it becomes unmaintainable.
  \item LaTeX debugging, table formatting, prose tightening.
  \item Adversarial pre-review of your own draft against the venue's form.
  \item Drafting rebuttal \emph{prose} from \emph{your} arguments.
  \item Turning notes into a first draft you then rewrite.
\end{itemize}
\end{multicols}

\begin{cbox}{A first review before any human reads it: Deep Review}
Google's \textbf{Deep Review} (\url{https://research.google/deepreview/}; public since
October 2026, after pilots on more than 16,000 manuscripts at STOC, ICML,
NeurIPS and ICLR) reads a PDF or plain text the way a careful referee would:
it splits the manuscript into sections, checks each in parallel for
mathematical errors, logical gaps and unsupported claims, then cross-checks
the whole. Two reviews a week are free; more need a Gemini API key. Google
states that uploads are not used to train its models.
\begin{itemize}
  \item Run it on the final PDF before every submission, then \emph{read every
        flag}: fix what is wrong, and write down why you reject what is not.
        Treat it as one more hostile reader, never as a verdict.
  \item It is for \emph{your own} drafts. Papers you are reviewing stay
        confidential (rule 3 below), and anything under NDA stays off the
        service whatever its training policy says.
\end{itemize}
\end{cbox}

\subsection{Non-negotiable guardrails}

\begin{wbox}{Five rules}
\begin{enumerate}
  \item \textbf{Never cite a paper you have not opened.} Models invent
        plausible references --- correct-looking authors, venue, year, all
        fabricated. Every entry in your \code{.bib} must come from Zotero,
        DBLP, or the Anthology, after you read at least the abstract.
  \item \textbf{Verify every number and every claim} against the source.
        Treat assistant output as a competent draft from a stranger, not as a
        result.
  \item \textbf{Never paste confidential data} --- NDA'd industry data,
        unpublished third-party manuscripts, papers you are reviewing --- into
        a cloud assistant. Reviewing material is confidential; using an LLM to
        write reviews is banned at most venues.
  \item \textbf{Know your venue's disclosure policy} and follow it. You are
        the author and you are responsible for every word.
  \item \textbf{Keep the core yours.} Derive your own proofs, design your own
        experiments, form your own judgements. The point of a PhD is that
        \emph{you} become the expert; outsource the mechanics, never the
        thinking.
\end{enumerate}
\end{wbox}

%==============================================================================
\section{Your personal site: \texttt{username.github.io}}
\label{sec:website}

You need one web page that you own, that outlives every institution you will
work at, and that answers ``who is this person and what have they done?'' in
ten seconds. GitHub Pages makes it free and permanent-ish; a Jekyll template
makes it take an afternoon.

\begin{cbox}{Why it matters more than it looks}
Programme chairs, reviewers who found your paper, internship recruiters,
postdoc hosts, and students looking for a supervisor all search your name. If
the first hit is a stale institutional directory page, you lose the
introduction. Your site is also where your PDFs, code links, slides, and CV
actually stay reachable after you change jobs.
\end{cbox}

\subsection{How GitHub Pages works}

Create a repository named \textbf{exactly} \code{username.github.io} (your own
GitHub username). Its contents are served at
\code{https://username.github.io}. One such \emph{user site} per account; free;
HTTPS included; custom domains supported. Jekyll is built in, so a template of
Markdown files becomes a site with no build pipeline of your own.

\subsection{Use the AcademicPages template}

\textbf{AcademicPages} --- \url{https://github.com/academicpages/academicpages.github.io}
--- is the standard academic Jekyll template (built on Minimal Mistakes). It
ships exactly the pages an ML/NLP researcher needs: publications, talks,
teaching, CV, blog posts, and a portfolio.

\begin{enumerate}
  \item On the AcademicPages repository, click \textbf{Use this template}
        $\rightarrow$ \textbf{Create a new repository}; name it
        \code{username.github.io}. (Cleaner than forking, which complicates
        pull requests and search.)
  \item Edit \code{\_config.yml}: your name, title, email, bio, location,
        GitHub, Scholar, ORCID, LinkedIn, and the site URL. This file alone
        controls most of what visitors see.
  \item Replace the content, which is all Markdown with YAML front matter:
        \begin{itemize}
          \item \code{\_pages/about.md} --- the landing page. Three paragraphs
                maximum: who you are, what you work on, what you are looking
                for.
          \item \code{\_publications/} --- one file per paper (title, venue,
                date, links to PDF / arXiv / code / BibTeX).
          \item \code{\_talks/}, \code{\_teaching/}, \code{\_posts/},
                \code{\_portfolio/} --- same pattern.
          \item \code{files/} --- your CV PDF and paper PDFs;
                \code{images/} --- your profile photo.
        \end{itemize}
  \item Generating publications from BibTeX: the scripts in
        \code{markdown\_generator/} turn a \code{.bib} or TSV export into the
        per-publication Markdown files. Export from Zotero, run the script,
        commit --- your site then updates from the same library that feeds
        your papers.
  \item Commit and push. GitHub builds and publishes within a minute or two;
        check the Actions tab if it does not appear. Optional local preview:
        \code{bundle install \&\& bundle exec jekyll serve}.
  \item Optional: a custom domain --- add a \code{CNAME} file and point your
        DNS at GitHub Pages. Worth it if you own \code{firstname-lastname.com},
        but \code{username.github.io} is perfectly respectable.
\end{enumerate}

\begin{cbox}{A worked example: \url{https://y3nk0.github.io}}
My own site is built on AcademicPages. Copy its structure rather than its
content:

\begin{itemize}
  \item \textbf{About} --- short bio with research interests (NLP, ML,
        robustness, reasoning, interpretability, trustworthy models), plus
        background and affiliations.
  \item \textbf{News} --- a running list of accepted papers, talks, and grants.
        The cheapest way to look active, and the first thing returning
        visitors read.
  \item \textbf{Research} --- themes rather than a paper dump, so a visitor
        learns what you \emph{work on}, not just what you published.
  \item \textbf{Publications} --- venue, year, and links out to PDF and code.
  \item \textbf{Talks} and \textbf{Teaching} --- separate pages, which matter
        for academic job applications.
  \item \textbf{CV} --- both a downloadable PDF and an HTML version.
\end{itemize}
Because AcademicPages sites are public repositories, you can inspect the
source of any site you like --- mine included
(\code{github.com/y3nk0/y3nk0.github.io}) --- and borrow its structure
directly.
\end{cbox}

\begin{itemize}
  \item \textbf{Alternatives:} \textbf{al-folio} (Jekyll, very popular in ML,
        BibTeX-driven publications page), \textbf{Minimal Mistakes} (the base
        theme, if you want full control), \textbf{Hugo Blox} (formerly
        Wowchemy), or a \textbf{Quarto} website if you already write in Quarto.
        Any of them beats having no site.
  \item \textbf{Keep it current.} Update after every acceptance, every talk,
        and every CV change. A site last touched three years ago reads as
        worse than no site.
\end{itemize}

%==============================================================================
\section{Research infrastructure and habits}
\label{sec:infra}

The tools above handle literature and writing. These handle the part that
actually consumes your time: experiments.

\subsection{Make experiments reproducible from day one}

\begin{itemize}
  \item \textbf{One git repository per project.} Code, configs, notes,
        the paper. Never ``final\_v3\_real.py''.
  \item \textbf{Deterministic environments:} \code{uv} or conda with a
        lockfile, or Docker/Apptainer if your cluster supports it. Record
        versions of everything, including CUDA.
  \item \textbf{Config files, not command-line archaeology} (Hydra,
        or plain YAML). Every run must be reconstructible from a committed
        config plus a commit hash.
  \item \textbf{Seeds and variance.} Report mean $\pm$ std over at least 3
        seeds. A single-seed improvement of 0.3 points is noise, and a
        reviewer will say so --- correctly.
  \item \textbf{Experiment tracking} from your first run: Weights \& Biases,
        MLflow, or TensorBoard. Log configs, metrics, and the commit hash.
        ``Which run produced Table 3?'' must have an answer in month nine.
  \item \textbf{A results file, not screenshots.} Your tables and figures
        should regenerate from logged results with one script.
  \item \textbf{Compute:} learn your lab's cluster and Slurm early (job arrays,
        checkpoint-and-resume, time limits); know what free or national
        resources you can apply for; and treat ``can I actually run this?'' as
        a design constraint when choosing problems.
\end{itemize}

\subsection{Habits that compound}

\begin{description}
  \item[Weekly advisor meeting with a written agenda.] Send 5 bullets
        beforehand: what I did, what I found, what is blocked, what I propose,
        what I need from you. Keep a shared running document. This one habit
        separates students who progress from students who drift.
  \item[A dated lab notebook.] Including negative results --- especially
        negative results, since you will otherwise repeat them.
  \item[Read daily, 20 minutes.] Titles and abstracts from the arXiv feed,
        then one paper properly. Join or start a reading group; explaining a
        paper out loud exposes exactly what you did not understand.
  \item[Write continuously.] Keep a living draft of the current paper from
        week one; writing is how you discover that your story has a hole.
  \item[Say yes to reviewing] (workshops first, then ARR) once you have a
        paper or two. Nothing teaches you what makes a paper acceptable faster
        than reviewing eight of them.
  \item[Be present in the community.] Follow researchers in your area, go to
        your local ML/NLP meetups, apply to summer schools and doctoral
        consortia, volunteer at a conference you cannot otherwise afford, and
        talk to people at posters --- including, especially, posters you only
        half understand.
  \item[Protect your mind.] Impostor feelings are near-universal and peak in
        year one; rejection is the default outcome of submission, not a verdict
        on you. Keep interests outside the PhD, keep a list of things that have
        worked, and talk to people when it gets heavy.
\end{description}

%==============================================================================
\section{Checklists}

\subsection*{Week 1}
\begin{multicols}{2}
\begin{itemize}
  \item Zotero + connector + Better BibTeX installed; auto-export configured.
  \item Google Scholar profile created; 3 alerts live.
  \item ORCID, arXiv, OpenReview, GitHub, Hugging Face accounts.
  \item Overleaf account; institutional licence checked.
  \item Deadline tracker bookmarked; target venues in the calendar.
  \item AI assistant set up with a project memory file.
  \item Project repository created with \code{notes/} and a README.
  \item Cluster access working; a trivial job submitted successfully.
\end{itemize}
\end{multicols}

\subsection*{Month 1}
\begin{itemize}
  \item Survey question agreed in writing with your advisor; criteria written
        down.
  \item Seed papers read with the three-pass method; extraction sheet started.
  \item Search log running (queries, dates, hit counts).
  \item Weekly advisor meeting slot fixed, with the 5-bullet agenda format.
  \item Reading group joined.
\end{itemize}

\subsection*{Months 2--3}
\begin{itemize}
  \item 60--120 papers screened and extracted; taxonomy drawn.
  \item 10--15 pages of survey written; three candidate gaps written as
        testable questions.
  \item One baseline reproduced end to end, tracked, with numbers you trust.
  \item \code{username.github.io} live, with About, Publications (even if
        empty), CV.
  \item A target venue chosen and backward-planned from its deadline.
\end{itemize}

\subsection*{Before every submission}
\begin{multicols}{2}
\begin{itemize}
  \item Every DOI and arXiv id in \code{refs.bib} resolves; titles match;
        nothing cited that you have not opened.
  \item Published versions cited where they exist; \code{bibtex-tidy} run.
  \item Capitals protected in titles; no ``arXiv preprint'' for published
        work.
  \item \code{arxiv-latex-cleaner} run; no names in comments or PDF
        metadata; code anonymised if the venue is double-blind.
  \item Limitations section and the venue's checklist completed.
  \item The venue's AI-assistance disclosure made, if required.
  \item A Deep Review pass on the final PDF; every flag fixed or consciously
        rejected.
  \item Every table and figure regenerated by script from logged results;
        seeds and variance reported.
  \item Page limit, template and anonymity window checked on the venue
        page, not the aggregator.
\end{itemize}
\end{multicols}

\subsection*{Year 1}
\begin{itemize}
  \item One submission out --- workshop or main conference. Submitted beats
        perfect.
  \item A rebuttal written (your first one will teach you more than the paper
        did).
  \item A talk given somewhere: lab seminar, reading group, local meetup.
  \item Code and checkpoints released for at least one project.
  \item A one-page research statement you can hand to anyone who asks what you
        work on.
\end{itemize}

%==============================================================================
\appendix
\section{Link directory}

\begin{longtable}{L{4.3cm} L{10.5cm}}
\toprule
\thead{What} & \thead{Where} \\
\midrule
\endfirsthead
\toprule
\thead{What} & \thead{Where} \\
\midrule
\endhead
\multicolumn{2}{l}{\textbf{\color{brand}Deadlines and venues}}\\
\textbf{PwC AI deadlines} & \url{https://paperswithcode.co/ai-deadlines} \\
\textbf{AI deadlines (ML/CV/NLP)} & \url{http://aideadlines.org/?sub=ML,CV,NLP} \\
AI deadlines (mirror) & \url{https://huggingface.co/spaces/huggingface/ai-deadlines} \\
Deadline source repo & \url{https://github.com/huggingface/ai-deadlines} \\
CCF deadlines & \url{https://ccfddl.com} \\
CORE conference ranks & \url{https://portal.core.edu.au/conf-ranks/} \\
CSRankings & \url{https://csrankings.org} \\
Scholar venue metrics & \url{https://scholar.google.com/citations?view_op=top_venues} \\
\midrule
\multicolumn{2}{l}{\textbf{\color{brand}Writing and sharing}}\\
Overleaf & \url{https://www.overleaf.com} \\
Prism (closes 1 Jan 2027) & \url{https://prism.openai.com} \\
Deep Review (Google) & \url{https://research.google/deepreview/} \\
arXiv & \url{https://arxiv.org} \\
arxiv-latex-cleaner & \url{https://github.com/google-research/arxiv-latex-cleaner} \\
Anonymous code sharing & \url{https://anonymous.4open.science} \\
ORCID & \url{https://orcid.org} \\
latexdiff & \url{https://ctan.org/pkg/latexdiff} \\
CITATION.cff & \url{https://citation-file-format.github.io} \\
Zenodo & \url{https://zenodo.org} \\
\midrule
\multicolumn{2}{l}{\textbf{\color{brand}References and search}}\\
Zotero & \url{https://www.zotero.org} \\
Better BibTeX & \url{https://retorque.re/zotero-better-bibtex/} \\
bibtex-tidy & \url{https://github.com/FlamingTempura/bibtex-tidy} \\
Google Scholar & \url{https://scholar.google.com} \\
Scholar Labs & \url{https://scholar.google.com/scholar_labs/search} \\
Semantic Scholar & \url{https://www.semanticscholar.org} \\
OpenAlex & \url{https://openalex.org} \\
ACL Anthology & \url{https://aclanthology.org} \\
OpenReview & \url{https://openreview.net} \\
DBLP & \url{https://dblp.org} \\
Connected Papers & \url{https://www.connectedpapers.com} \\
ASReview & \url{https://asreview.nl} \\
Scite & \url{https://scite.ai} \\
Unpaywall & \url{https://unpaywall.org/products/extension} \\
HF Daily Papers & \url{https://huggingface.co/papers} \\
Papers with Code & \url{https://paperswithcode.co} \\
\midrule
\multicolumn{2}{l}{\textbf{\color{brand}Attention signals}}\\
Star history & \url{https://star-history.com} \\
Hugging Face hub & \url{https://huggingface.co/models} \\
\midrule
\multicolumn{2}{l}{\textbf{\color{brand}Personal site}}\\
AcademicPages template & \url{https://github.com/academicpages/academicpages.github.io} \\
GitHub Pages docs & \url{https://docs.github.com/pages} \\
Example site & \url{https://y3nk0.github.io} \\
al-folio & \url{https://github.com/alshedivat/al-folio} \\
\bottomrule
\end{longtable}

\section{A minimal paper repository}

\begin{quote}\small\ttfamily
project/\\
\hspace*{1em}README.md\\
\hspace*{1em}CITATION.cff\hfill \rmfamily\textit{citable repo; Zenodo DOI per release}\ttfamily\\
\hspace*{1em}CLAUDE.md\hfill \rmfamily\textit{conventions for your AI assistant}\ttfamily\\
\hspace*{1em}pyproject.toml / environment.yml\hfill \rmfamily\textit{pinned}\ttfamily\\
\hspace*{1em}configs/\hfill \rmfamily\textit{one YAML per experiment}\ttfamily\\
\hspace*{1em}src/\\
\hspace*{1em}scripts/\hfill \rmfamily\textit{train.sh, eval.sh, slurm/}\ttfamily\\
\hspace*{1em}notes/\hfill \rmfamily\textit{2026-10-15.md, dated lab notebook}\ttfamily\\
\hspace*{1em}results/\hfill \rmfamily\textit{logged metrics, committed}\ttfamily\\
\hspace*{1em}figures/make\_figures.py\hfill \rmfamily\textit{results $\rightarrow$ PDF}\ttfamily\\
\hspace*{1em}paper/\\
\hspace*{2em}main.tex, sections/, refs.bib\hfill \rmfamily\textit{BBT auto-export target}\ttfamily
\end{quote}

\vspace{0.6cm}
\begin{cbox}{The shortest version of all of this}
Keep a clean library, know your deadlines, read forward from the anchor papers,
reproduce before you innovate, write every week, publish a page with your name
on it, use the assistant for the mechanics and never for the thinking --- and
ask your advisor the question you are embarrassed to ask. Good luck.
\end{cbox}

\end{document}
